\scanpage{79}
\mspage{85}

Dém. a) Il suffit, en vertu de prop. 3, de montrer que l'hyp. implique que pour \uline{tout} $s\in S$, $f$ induit une $\mathbb D$-équivalence $X_s\xrightarrow{f_s}X'_s$. Cela résulte \deleted{\unsure{du le}} aussitôt du \deleted{lemme}

\uline{Lemme} \deleted{\unsure{sous les hyp}} Si $p$ et $p'$ sont tous deux \unsure{soient|soit} des $\mathbb D$-fibrations gauches, \unsure{soient|soit} des $\mathbb D$-fibr. droites, et si $s_1\longrightarrow s_0$ est une flèche dans $S$, alors $f_{s_1}$ est une $\mathbb D$-équiv. sss $f_{s_0}$ l'est.

(\deleted{Donc} \inserted{L'ens $S'\subset\mathrm{Ob}\,S$} des $s\in S$ tels que $f_s$ soit une $\mathbb D$-équiv., est une réunion de comp. connexes de $S$.)

Supposons p. ex. que $p,p'$ soient des $\mathbb D$-fibr. gauches. Considérons \inserted{pour $s\in S$} le diagramme
\diagram{\begin{tikzcd}[ampersand replacement=\&]
X_{s\deleted{\unsure{1}}}\arrow[r,"f_s"']\arrow[r,phantom,shift left=1.6ex,"{\deleted{\unsure{\longleftrightarrow|\longrightarrow}}}" description]\arrow[d] \& X'_{\unsure{s|s_0}}\arrow[d]\\
X_{/\unsure{s|s_0}}\arrow[r,"f_{/s}"] \& X'_{/\unsure{s|s_0}}
\end{tikzcd}}
Les flèches verticales sont des $\mathbb D$-équiv. par l'hyp. sur $p,p'$, donc $f_s$ est une \inserted{$\mathbb D$-}équiv. sss $f_{/s}$ l'est.

Considérons maintenant
\diagram{\begin{tikzcd}[ampersand replacement=\&]
X_{/s_1}\arrow[r,"f_{\unsure{s_1|/s_1}}"]\arrow[d] \& X'_{/s_1}\arrow[d]\\
X_{/s_0}\arrow[r,"f_{/s_0}"] \& X'_{/s_0}
\end{tikzcd}}
Ici encore, \unsure{les|l'} hyp sur $p,p'$ \unsure{impliquent|implique} que les flèches verticales sont des $\mathbb D$-équiv., donc $f_{s_1}$ est une $\mathbb D$-équiv. sss $f_{s_0}$ l'est.
% archive sheet number: 78
